\documentclass[11pt]{article}
\usepackage{mathblatt}
\begin{document}
\blattkopf*{Basisaufgaben}{BAS-Z2}{Basisaufgaben · Zettel · BAS-Z2}
\einheitenkopf[][Zettel 2]{Basisaufgaben · Zettel 2}
\zweigzeile{ohne Taschenrechner · je Aufgabe 1 Punkt · Lösungen auf der Rückseite}
% \small: zehn Aufgaben mit bis zu zwei Koordinatensystemen auf einer Seite (Render 28.09.)
\small

% 1: Bruchteil einer Fläche bestimmen – brueche-dezimalzahlen-basis-k1-v2 (Original 2026-FOR-B1b)
\begin{aufgabe}{Bei welcher Figur ist genau $\frac{2}{3}$ grau? \hfill (P26F) \\ \kreuz{P} \kreuz{Q} \kreuz{R} \kreuz{S}}

P \bruchrechteck[2.5]{2}{5} \quad Q \bruchkreis[0.8]{3}{5} \quad R \kreissektor[0.8]{240}{} \quad S \bruchkreis[0.8]{3}{4}
\end{aufgabe}

% 2: Zahlen in verschiedenen Darstellungen vergleichen – brueche-dezimalzahlen-basis-k4-v2 (Original 2023-OS-B1f)
\begin{aufgabe}{Welcher Wert ist der kleinste: $6{,}6$; $0{,}6^2$; $0{,}66$; $66\,\%$? \hfill (P23) \\ \leerfeld}
\end{aufgabe}

% 3: Zehnerpotenzschreibweise umwandeln – potenzen-wurzeln-basis-k3-v1 (Original 2025-OS-B1d)
\begin{aufgabe}{$62\,000 = 6{,}2 \cdot 10^{\square}$ – welche Hochzahl? \hfill (P25) \\ \leerfeld}
\end{aufgabe}

% 4: Grundwert berechnen – prozentrechnung-basis-k1-v1 (Original 2025-OS-B1a)
\begin{aufgabe}{Beim Kauf einer Jacke spart Tom $8\,€$, das sind $20\,\%$ des alten Preises. Wie hoch war der alte Preis? \hfill (P25) \\ \leerfeld[€]}
\end{aufgabe}

% 5: Prozent und Anteil umwandeln – prozentrechnung-basis-k2-v1 (Original 2022-OS-B1f)
\begin{aufgabe}{Jeder zehnte Besucher gewinnt einen Preis – wie viel Prozent? Kreuze an. \hfill (P22) \\ \kreuz{$1\,\%$} \kreuz{$10\,\%$} \kreuz{$20\,\%$} \kreuz{$90\,\%$}}
\end{aufgabe}

% 6: Term zu Sachtext angeben – terme-basis-k2-v1 (Original 2023-OS-B1h)
\begin{aufgabe}{Die Summe aus dem Doppelten einer Zahl $x$ und 5 wird verdreifacht. Welcher Term passt? \hfill (P23) \\ \kreuz{$(2x + 5) : 3$} \kreuz{$2x + 5 \cdot 3$} \kreuz{$3 \cdot (2x + 5)$} \kreuz{$2x \cdot 5 + 3$}}
\end{aufgabe}

% 7: Figur nach Spiegelung benennen – symmetrie-abbildungen-basis-k1-v1 (Original 2018-OS-B1h)
\begin{aufgabe}{Ein Dreieck mit drei verschieden langen Seiten liegt mit einer Seite auf der Geraden $g$ und wird an $g$ gespiegelt. Welches Viereck bilden Dreieck und Spiegelbild? \hfill (P18)}
\end{aufgabe}

% 8: Gleichschenkliges Dreieck erkennen – winkel-dreiecke-basis-k2-v1 (Original 2023-OS-B1g)
\begin{aufgabe}{Im Dreieck ABC ist $\alpha = \beta = 65^\circ$ und AC = $5$ cm. Gib die Länge von BC und die Größe von $\gamma$ an. \hfill (P23) \\ BC = \leerfeld[cm] , $\gamma =$ \leerfeld °}
\end{aufgabe}

% 9: Gegenfläche im Würfelnetz bestimmen – koerper-basis-k1-v1 (Original 2016-OS-B1j)
\begin{aufgabe}{\begin{minipage}[t]{0.6\linewidth}Das Netz aus sechs Quadraten (je Buchstabe ein Quadrat) wird zu einem Würfel gefaltet. Der Würfel wird auf die Fläche C gestellt. Welche Fläche liegt oben? \hfill (P16) \\ \leerfeld\end{minipage}\hfill\begin{minipage}[t]{0.37\linewidth}\vspace{0pt}
$\begin{array}{cccc} A &  &  &  \\ B & C & D & E \\ F &  &  &  \end{array}$
\end{minipage}}
\end{aufgabe}

% 10: Winkelfunktion Seitenverhältnis angeben – trigonometrie-basis-k2-v1 (Original 2025-OS-B1g)
\begin{aufgabe}{\begin{minipage}[t]{0.6\linewidth}Rechter Winkel bei $B$. Trage den Bruch ein. \hfill (P25) \\ $\mathrm{sin}\,\alpha =$ \leerfeld\end{minipage}\hfill\begin{minipage}[t]{0.37\linewidth}\vspace{0pt}
\dreieck{(0,0)}{(3,0)}{(3,2.2)}{k}{m}{n}{\alpha}{}{}
\end{minipage}}
\end{aufgabe}

\begleitteil
\erg{1}{R, denn $\frac{240}{360} = \frac{2}{3}$}
\erg{2}{$0{,}6^2 = 0{,}36$; die anderen sind $6{,}6$ und zweimal $0{,}66$}
\erg{3}{$4$ – das Komma wandert vier Stellen nach links}
\erg{4}{$8 \cdot 5 = 40\,€$ (nicht $20\,\%$ von $8\,€$)}
\erg{5}{$10\,\%$ ($\frac{1}{10} = \frac{10}{100}$)}
\erg{6}{$3 \cdot (2x + 5)$}
\erg{7}{Drachenviereck}
\erg{8}{BC = AC = $5$ cm (gleiche Winkel bei A und B); $\gamma = 180^\circ - 2 \cdot 65^\circ = 50^\circ$}
\erg{9}{E (E und C liegen sich gegenüber)}
\erg{10}{$\mathrm{sin}\,\alpha = \frac{k}{m}$}
\end{document}
